\textbf{Время:} резерв

\textbf{Тезис:} текущие радиусы исследованы в десяти соседних интервалах быстроты от −1 до 1.

\textbf{На экране:} три радиуса слева направо, горизонтальная ось y и общая легенда центральности.

\textbf{Текст:} Быстрота описывает продольное движение пары. Здесь она определена в системе центра масс Au+Au; рассмотрен диапазон от минус единицы до единицы с шагом 0,2. Каждая точка относится к отдельному интервалу, а цвета снова различают центральности.

Сравним панели: side меняется слабее, а в out и long для ряда центральностей заметна зависимость от положения по быстроте. Это ещё одна зависимость, с которой можно сравнивать сценарии продольного расширения. Саму форму кривых не следует отождествлять с прямым профилем плотности источника.

\textbf{Переход:} На следующем резервном слайде — параметр лямбда тех же гауссовых описаний.

\textbf{Справка:} c\_all\_graphs\_neg.jpg. Диапазон и шаг подтверждены автором. Для пары \(y=\tfrac12\ln[(E_{\rm pair}+P_{z,\rm pair})/(E_{\rm pair}-P_{z,\rm pair})]\). Система центра масс столкновения не является LCMS всех пар сразу: переход в LCMS выполняется для каждой пары. Условия усреднения по kT в предоставленном файле не записаны. Кривые не симметризованы заново, исходные точки сохранены.

\textbf{Источники:} авторский JPEG от 19.10.2026 и текущий манифест.
